% FORMLÆRE - Ei traktatform om korleis form oppstår
% Iver Raknes Finne, AHO 2026
% Spraak: Nynorsk

\documentclass[11pt, a4paper]{article}

\usepackage[utf8]{inputenc}
\usepackage[T1]{fontenc}
\usepackage[nynorsk]{babel}
\usepackage{mathpazo}
\usepackage{amsmath, amssymb}
\usepackage{hyperref}
\usepackage{setspace}
\usepackage{geometry}
\usepackage{csquotes}
\usepackage{enumitem}

\geometry{margin=2.5cm, top=2.5cm, bottom=2.5cm}
\setstretch{1.15}

\newcommand{\prop}[1]{\bigskip\noindent{\Large\textbf{#1}}\medskip\par}
\newcommand{\subprop}[1]{\medskip\noindent\textbf{#1}\quad}
\newcommand{\subsubprop}[1]{\smallskip\hangindent=1.5em\noindent\hspace{1.5em}\textbf{#1}\quad}

\title{%
  \vspace{-1cm}
  {\huge\textbf{FORML\AE RE}}\\[0.5em]
  {\large\textit{Ei traktatform om korleis form oppst\aa r}}
}
\author{
  iver.raknes.finne@stud.aho.no\\
  Arkitektur- og designh\o gskolen i Oslo\\
  2026
}
\date{}

\begin{document}
\maketitle
\thispagestyle{empty}

\section*{Forord}

Denne teksten er skriven som ein traktat: ei rekkje nummererte proposisjonar der kvar p\aa stand utdjupar, avgrensar eller dreg konsekvensar av dei f\o reg\aa ande. Ho gjer krav p\aa\ \aa\ vere fullstendig i sin logikk, ikkje i sine d\o me.

Rammeverket er meint som eit teoretisk verkt\o y for alle felt der form vert gjeven til materie, og for dei felta der form oppst\aa r utan at nokon gjev ho. Det gjev grunnlag for kvantitativ analyse og forklaringsmodellar av stil og periode, for forst\aa ing av generativ og parametrisk formgjeving, og for samanstilling av naturleg og kunstig form under same logiske struktur.

Proposisjonane er pr\o vde mot eit datasett p\aa\ om lag 2300 europeiske stolar fr\aa\ Nasjonalmuseet og Victoria and Albert Museum, som spenner fr\aa\ 1280 til 2024. Kvar generell p\aa stand har s\aa leis ein empirisk ber\o ringsflate i stolformene. Men teksten er ikkje avgrensa til stolar: ho gjeld kvar form som oppst\aa r under avgrensande vilk\aa r.

\newpage
\tableofcontents
\newpage

% ============================================
\prop{1 Ting har former.}

\subprop{1.1} Alt som finst i rommet, har ei form. Ei h\o gd, ei breidd, ei djupn, eit materiale, ein proporsjon, ei overflate. Summen av alt dette er forma.

\subprop{1.2} Totaliteten av alle former ein ting kunne ha hatt, er formrommet hans.

\subsubprop{1.21} Formrommet er ikkje ein abstraksjon. Det er den reelle mengda av alle moglege kombinasjonar av geometri, materiale og konstruksjon. Kvar dimensjon svarar til ein eigenskap som kan variere. Kvar realisert gjenstand er eitt punkt i dette rommet.

\subsubprop{1.22} Formrommet har tre typar regionar: busette (former som faktisk finst), tomme men moglege (former som ikkje bryt fysiske lovar, men som ingen har laga), og forbodne (former som ikkje kan eksistere). Grensa mellom dei to siste flyttar seg med teknologien.

\subsubprop{1.23} Ikkje alle posisjonar er busette. Dei tome plassane er like informative som dei fulle: dei fortel kva former som ikkje var moglege, eller ikkje overlevde.

\subsubprop{1.24} Dei busette regionane utgjer ein forsvinnande liten del av formrommet. I kvart system som er unders\o kt, fr\aa\ biologiske skjelformer til datalagde stolmodellar, okkuperer dei realiserte formene under \'ein prosent av det teoretisk moglege rommet.

\subsubprop{1.25} Formrommet kan m\aa last. Kvar dimensjon svarar til ein m\aa lbar eigenskap: h\o gd, breidd, djupn, vekt, materialtype, samanf\o yingsmetode. Tettleiken av realiserte former i ein region er eit m\aa l p\aa\ kor stabil regionen er.

\subprop{1.3} Sidan formrommet inneheld fleire busette posisjonar (1.23), er det eit forklaringskrevjande faktum at ein gjenstand okkuperer \'ein posisjon og ikkje ei anna.

\subprop{1.4} All formgjeving gjeld dette grunnfaktumet. Forklaringa av det er emnet for alt som f\o lgjer.

% ============================================
\prop{2 Forma er bestemt av mange krefter samstundes.}

\subprop{2.1} Kvar ting som vert til, vert til under eit sett av vilk\aa r. Vilk\aa ra er aldri eitt; dei er fleire, og dei verkar p\aa\ same materie p\aa\ same tid.

\subsubprop{2.11} Kva stoff som er tilgjengeleg. 2.12 Kva reiskapar og teknikkar som finst. 2.13 Kva det kostar. 2.14 Kvar i verda tingen vert til. 2.15 Kva omgjevnaden ventar, premierer og tolererer. 2.16 Kva kroppen bruk krev. 2.17 Kva regulatoriske system som avgrensar. 2.18 Kva tidlegare former som kan kopierast eller avvisast.

\subsubprop{2.19} Kvar av desse er eit \textit{seleksjonstrykk}: ein faktor som endrar sannsynlegheitsfordelinga over formrommet. Funksjonen (2.16) er eitt slikt trykk. Han er reell. Men han er berre eitt blant fleire.

\subprop{2.2} Om eitt einaste seleksjonstrykk avgjorde forma, ville alle ting underlagde same trykk hatt same form.

\subsubprop{2.21} Dersom funksjonen \aa leine bestemte forma, ville alle ting med same funksjon sett like ut.

\subsubprop{2.22} Sj\aa\ p\aa\ kva stol som helst: dei har alle same funksjon, men radikalt ulike former. Dette er eit empirisk faktum. Det stadfestar at fleire uavhengige seleksjonstrykk enn funksjon verkar.

\subsubprop{2.23} Innanfor klassen ``stol'' varierer volumet med ein faktor p\aa\ fleire hundre. Same funksjon, radikalt ulik utstrekning. H\o gda ber over eitt bit gjensidig informasjon med stilperioden. Funksjonen, som er konstant, ber null.

\subsubprop{2.24} Funksjonen avgrensar formrommet. Ho eliminerer dei forbodne regionane. Men innanfor dei gjenverande regionane er funksjonen taus: ho forklarer kvifor visse former er umoglege, men ikkje kvifor \'ei bestemt form vart realisert framfor ei anna.

\subprop{2.3} Seleksjonstrykka dreg sjeldan i same retning.

\subsubprop{2.31} Det materialet tilbyr, strider mot det \o konomien krev. Det kulturen ventar, strider mot det geografien t\aa ler. Det kroppen treng, strider mot det teknologien meistrar.

\subprop{2.4} Sidan fleire uavhengige trykk verkar samstundes (2.22), og sidan dei dreg i ulike retningar (2.3), kan ikkje alle verte fullt ut tilfredsstilte.

\subsubprop{2.41} Kvar realisert form er difor eit kompromiss. Ikkje eit d\aa rleg kompromiss, men det einaste moglege kompromisset under dei vilk\aa ra som r\aa dde.

\subprop{2.5} Kvart seleksjonstrykk ber ein m\aa lbar mengd informasjon om den endelege forma. Ein maskinl\ae ringsmodell som brukar materiale, tid og produksjonsland oppn\aa r vesentleg betre treffsikkerheit enn ein som berre brukar funksjon. Forklaringskrafta ligg i kombinasjonen, ikkje i einskildfaktorane.

\subsubprop{2.51} Formelen ``form f\o lgjer funksjon'' er ikkje feil. Ho er eit grensetilfelle: ho gjeld berre der eitt einaste seleksjonstrykk dominerer s\aa\ sterkt at alle andre er negligerbare. I den reelle verda er dette unntaket, ikkje regelen.

% ============================================
\prop{3 Kreftene lagar eit tilpassingslandskap.}

\subprop{3.1} Seleksjonstrykka verkar ikkje kvar for seg; dei definerer saman eit \textit{tilpassingslandskap} over formrommet. Kvar posisjon i formrommet har ein verdi: kor godt forma tilfredsstiller alle verksame trykk samstundes.

\subsubprop{3.11} Landskapet er ikkje flatt. Det har haugar (former som er godt tilpassa), dalar (former som er d\aa rleg tilpassa) og sadelpunkt (former som er stabile i nokre retningar og ustabile i andre).

\subsubprop{3.12} Landskapet er ikkje glatt. Det har mange lokale haugar. Ein form kan vere den beste i sitt nabolag utan \aa\ vere den beste overalt.

\subsubprop{3.13} Talet p\aa\ haugar i landskapet aukar med talet p\aa\ samverkande seleksjonstrykk. Eitt trykk \aa leine gjev \'ein haug. Fleire uavhengige trykk som verkar samstundes gjev mange haugar. Jo fleire gjensidig avhengige trykk, jo meir ulendt vert landskapet.

\subprop{3.2} Kvar haug svarar til eit \textit{arketype}: ei form som er stabil fordi sm\aa\ endringar i alle retningar gjev d\aa rlegare tilpassing.

\subsubprop{3.21} Arketypar er ikkje oppfinningar. Dei er attraktorar i landskapet. Dei ``finst'' der uavhengig av kven som oppdagar dei.

\subsubprop{3.22} Ein \textit{stil} er ei klynge av realiserte former rundt same haug. Spreiinga i klynga er eit m\aa l p\aa\ kor flat haugen er.

\subsubprop{3.23} Eit arketype som dukkar opp uavhengig i tradisjonar utan kontakt, er sterkt prov for at haugen er ein eigenskap ved landskapet, ikkje ved nokon einskild vandrar.

\subprop{3.3} Formgjeving er ikkje skapande i den forstand at navigatoren konstruerer landskapet. Landskapet er gjeve av vilk\aa ra. Navigatoren vandrar i det.

\subsubprop{3.31} \AA\ oppdage ein haug er ikkje \aa\ skape han. Det er \aa\ finne det landskapet allereie inneheld.

\subprop{3.4} Landskapet har ikkje \'ein global haug. Det har fleire lokale haugar, separerte av dalar.

\subsubprop{3.41} Kvar haug trekkjer til seg former i sitt nabolag. N\aa r eitt seleksjonstrykk dominerer, smeltar fleire haugar saman til \'ein. Landskapet vert glattare, og mangfaldet krympar.

\subsubprop{3.42} N\aa r fleire trykk er omtrent like sterke, held haugane seg \aa tskilde. Mangfaldet aukar.

% ============================================
\prop{4 Landskapet er i r\o rsle.}

\subprop{4.1} Seleksjonstrykka er ikkje stasjon\ae re. Dei endrar seg over tid.

\subsubprop{4.11} Ein ny teknologi opnar ein region av formrommet som var forboden (1.22). Ein ny marknad endrar kva som vert premieret. Eit nytt materiale endrar kva som er mogleg.

\subprop{4.2} N\aa r seleksjonstrykka endrar seg, endrar landskapet seg. Haugar kan stige, senke seg, flytte seg, dele seg eller smelte saman.

\subsubprop{4.21} Ein form som var optimalt tilpassa under eitt sett av vilk\aa r, kan vere d\aa rleg tilpassa under eit anna.

\subprop{4.3} R\o rsla i landskapet er ikkje tilfeldig. Ho f\o lgjer strukturelle endringar i teknologi, \o konomi, demografi og kultur.

\subsubprop{4.31} Etterkvar omvelting opnar rommet seg att. Nye materialar, nye teknikkar, nye vilk\aa r gjev nye haugar. Over lang tid vert landskapet rikare.

\subsubprop{4.32} Nye haugar oppst\aa r ikkje berre ved at gamle flatar ut. Dei oppst\aa r \o g ved rekombinasjon: samankopling av element fr\aa\ ulike tradisjonar, materialar eller teknikkar som ikkje tidlegare har m\o tst. Rekombinasjon opnar regionar av formrommet som var utilgjengelege fr\aa\ kvar enkelt utgangstradisjon.

\subprop{4.4} Ein \textit{stilperiode} er eit augeblinksbilete av landskapet. Han fangar dei haugane som var aktive, dei regionane som var tilgjengelege, og dei kreftene som r\aa dde.

\subprop{4.5} Kvart realisert objekt er eit svar p\aa\ landskapet slik det var i det augeblinket objektet vart til. Det er gyldig for dei vilk\aa ra som r\aa dde d\aa\ det vart til. Vilk\aa ra er allereie i endring i det augeblinket objektet er ferdig.

\subprop{4.6} Landskapet har minne. Kvar realisert form etterlet seg eit spor: eit fysisk objekt, ein produksjonsmetode, ei forventning, ein vane. Sporet vert sj\o lv ein del av seleksjonstrykka for neste form.

\subprop{4.7} Endringa i landskapet er ikkje jamn. Lange periodar med sm\aa\ justeringar vert avbrotne av br\aa\ opningar: ein ny teknologi gjer tilgjengeleg ein heil region som f\o r var forboden. M\o nsteret er eit av stase og brot.

\subsubprop{4.71} Etter kvart brot f\o lgjer ei rask spreiing av nye former inn i den nyopna regionen, deretter ei gradvis innsnevring mot nye haugar. Dampb\o ying opna \'ein region; r\o yrst\aa l ein annan; spr\o ytest\o yping ein tredje. Kvar gong var m\o nsteret det same.

% ============================================
\prop{5 Materialet deltek i \aa\ avgjere forma.}

\subprop{5.1} Den r\aa dande modellen i formfaga har vore at formgjevaren tenkjer, og materialet lystrer. At form vert p\aa f\o rt materie utanfr\aa: ein plan fr\aa\ ein hjerne, utf\o rt p\aa\ eit passivt stoff.

\subsubprop{5.11} Denne modellen er utilstrekkeleg. Ho kan ikkje forklare kvifor same formgevar, med same intensjon og same funksjonskrav, produserer radikalt ulike former i ulike materialar.

\subsubprop{5.12} Skiljet mellom den som tenkjer og det som vert tenkt p\aa, er ikkje skarpt. Det er graduelt. Det finst ein samanhengande skala fr\aa\ system som berre kan endrast ved \aa\ omkalfatre maskinvara, via system der parametrane kan justerast, til system som justerer seg sj\o lve.

% ============================================
\prop{6 \AA\ gi form er \aa\ navigere.}

% ============================================
\prop{7 Navigasjon er substrat-uavhengig.}

% ============================================
\prop{8 Navigasjonen er alltid distribuert.}

% ============================================
\prop{9 Lova er stasjon\ae r verknad.}

% ============================================
\prop{10 Ingen form er endeleg.}

\end{document}
